\section{Day 1: Gaussian Integers}
The Quercus site for this class will be up tonight or early tomorrow. This class will have (approximately) weekly homework. There will be regular oral quizzes with the TA, about 5--10 minutes, where you will discuss one of the problems on the homework.

Consider the following example,
\[ 1+1 = 2, \quad 1+4 = 5, \quad 4+9 = 13, \quad 1+16=17. \]
\vspace{-22pt}
\begin{theorem}[Fermat] \label{thm:fermat_on_two_squares}
    Let $p > 2$; $p$ can be expressed as a sum of two squares if and only if $p \equiv 1 \pmod{4}$.
\end{theorem}
\begin{proof}
    We start with the easy direction. Observe that the squares of the numbers $0, 1, 2, 3, \dots$ are all equal to $0$ or $1$ modulo $4$, i.e., $0^2 = 0$, $1^2 = 1$, $2^2 = 4 \equiv 0 \! \pmod{4}$, and $3^2 = 9 \equiv 1 \! \pmod{4}$, and so on (you can always express a number as either $2n$ or $2n + 1$, then squaring, to check that this pattern continues). It follows that if $p = n^2 + m^2$, then $p \equiv 1 \! \pmod{4}$; if it were the sum of the squares of two odd numbers, then $p$ would be even, which contradicts it being prime.

    For the other direction, consider the ring of Gaussian integers,
    \[ \ZZ[i] = \{a + bi \mid a, b \in \ZZ\} \subset \CC. \]
    Consider $p = n^2 + m^2$. Notice that the right hand side cannot be factored as \textit{integers}, but it could be factored in $\ZZ[i]$ as $(n + im)(n - im)$; in this manner, we may instead ask the question of how to factor $p \in \ZZ[i]$. The continuation of this proof will use the fact that $\ZZ[i]$ is a UFD, but we will not prove that today. In particular, $\ZZ[i]$ is a Euclidean domain with $\alpha \mapsto \abs{\alpha}^2$ is a Euclidean norm.
\end{proof}

Given $z = x + iy \in \ZZ[i]$, define the norm of $z$ as $N(z) = z \cdot \ol{z} = x^2 + y^2 \in \ZZ_{\geq 0}$.\footnote{at this point, Rozenblyum made a comment about complex number norms versus algebra norms. idk. sure.} Observe that we have $N(z \cdot w) = N(z) \cdot N(w)$.
\begin{claim}
    $z \in \ZZ[i]$ is a unit if and only if $N(z) = 1$, i.e., the units are $\{\pm 1, \pm i\}$.
\end{claim}
\begin{proof}
    Suppose $N(z) = 1$; by definition, we see $z \cdot \ol z = 1$, and so $\ol{z}$ constitutes the inverse of $z$. Conversely, suppose $z \in \ZZ[i]$ is a unit; if there exists another number $w$ such that $z \cdot w = 1$, then, since norm is multiplicative, we may take the norm of both sides to get that $N(z) \cdot N(w) = 1$. Since $N(z) \in \ZZ_{\geq 0}$, it must be that $N(z) = 1$.
\end{proof}

\begin{claim} \label{clm:primes_in_gaussian_integers}
    If $p$ is a prime such that $p \equiv 1 \pmod{4}$, then $p$ is not prime in $\ZZ[i]$.
\end{claim}
We will prove the claim later; for now, observe that \ref{thm:fermat_on_two_squares} and \ref{clm:primes_in_gaussian_integers} imply each other; one is obvious, and for the direction that the claim implies the theorem, suppose $p$ is not prime, so we may write it as the product of two non-units, $p = \alpha \cdot \beta$, where $\alpha, \beta \in \ZZ[i]$. We have that $p^2 = N(p) = N(\alpha) \cdot N(\beta)$; but $N(\alpha), N(\beta) \neq 1$, so it must be that $N(\alpha) = N(\beta) = p$. By writing $\alpha = a + bi$, we get that $N(\alpha) = a^2 + b^2 = p$, which is the desired expression as seen in \ref{thm:fermat_on_two_squares}.

We now prove \ref{clm:primes_in_gaussian_integers} in two ways. For the first way, consider Wilson's theorem:
\begin{theorem}[Wilson] \label{thm:wilson}
    Let $p$ be prime; then $(p-1)! \equiv -1 \pmod{p}$.
\end{theorem}
\begin{proof}
    The proof of this is simple; since $\ZZ/p$ is a field, every nonzero residue modulo $p$ has an inverse, where $\alpha \equiv \alpha\inv$ if and only if $\alpha \equiv 1$ or $-1$, i.e., $\alpha^2 = 1$. It follows that the product of $1, \dots, p-1$ without $1$ and $p-1$ is $1$, so we are done.
\end{proof}
\begin{corollary} \label{cor:wilson}
    If $p = 4n + 1$, then $(\!(2n)!)^2 \equiv -1 \pmod{p}$.
\end{corollary}
\begin{proof}
    By Wilson's theorem, we have that $-1 \equiv (p-1)!$. Let us write the factorial out;
    \[ -1 \equiv (p-1)! = (1 \cdot 2 \cdot \ldots \cdot 2n) \bigl((p - 2n) (p - 2n + 1) \dots (p-1)\bigr). \]
    Since we're working modulo $p$, we may realize the right hand side expression as
    \[ \dots \equiv (1 \cdot 2 \cdot \ldots \cdot 2n)\bigl((-2n) (-2n + 1) \dots (-1)\bigr) = (2n)! (-1)^{2n} (2n)! = \bigl((2n)!\bigr)^2. \qedhere \]
\end{proof} \vspace{-8pt}
The above corollary allows us to write\footnote{Rozenblyum proceeded to stand stunned in front of the blackboard for 3 minutes. When pointed out on how to continue the discussion, he said that this was proof your ``IQ drops precipitously when in front of a classroom''.}
\[ p \mid \bigl((2n)!\bigr)^2 + 1, \]
which we may factor as $(\!(2n)! + i)(\!(2n)! - i)$; if $p$ was prime, then it divides one of them, but we know that $p (a + bi) = (2n)! \pm i$ can't happen.

We now proceed to the second proof. Consider the cyclic group
\[ (\ZZ/p)^\times = \ZZ/(p-1); \]
if $p = 4n + 1$, then $\ZZ/(p-1) \cong \ZZ/(4n)$, whence the existence of some element of order $4$ in $(\ZZ/p)^\times$. Let such an element be called $m$; it follows that $m^4 \equiv 1 \pmod{p}$ and $m^2 \equiv -1 \pmod{p}$, but then we would have $p \mid m^2 + 1$, which is a contradiction. While this proof was fast, it didn't really tell us much about the arithmetic that went on in the first proof.

The morals as to the two above proofs are that
\begin{enumerate}[(i)]
    \item It is important to understand arithmetic in more general number rings (such as prime factorization). As an example,
    \[ \ZZ[\sqrt{-5}] = \{a + b \sqrt{-5} \mid a, b \in \ZZ\} \]
    is not a UFD (see the MAT347 notes for a reference on this).
    \item Questions about number rings are closely related to questions in modular arithmetic.
\end{enumerate}
